%%%%%%%%%%  scan idx 29  |  folio 16  %%%%%%%%%%
\SGAPrintPage{16}
\[
f:X\longrightarrow Y
\]
is a proper morphism of noetherian schemes, then
\sgaStableTarget{sga6.e00.eq.6-2}\begin{equation}\tag{6.2}
f_*:K_\bullet(X)\longrightarrow K_\bullet(Y)
\end{equation}
is compatible with the filtrations and defines a homomorphism of graded groups
\sgaStableTarget{sga6.e00.eq.6-3}\begin{equation}\tag{6.3}
f_*:\Gr_\bullet(X)\longrightarrow \Gr_\bullet(Y),
\end{equation}
giving a ``projection formula'', i.e. it is $\Gr^\bullet(Y)$-linear, a formula which is deduced by passage to the associated graded objects from the analogous projection formula for the homomorphism (6.2) (which is $K^\bullet(Y)$-linear).

\uline{From the point of view proposed in the Seminar}, \uline{a workable theory of intersections}, \uline{on noetherian schemes} $X$ \uline{that are not necessarily regular}, \uline{must consist in the combined study of the contravariant invariants} $K^\bullet(X),\Gr^\bullet(X)$ \uline{and the covariant invariants} $K_\bullet(X),\Gr_\bullet(X)$, \uline{which play roles of an essentially different nature}, \uline{and whose properties complement one another}. Thus the covariant groups $K_\bullet(X),\Gr_\bullet(X)$ can be calculated more readily for certain nonproper fibrations, by means of the ``homotopy exact sequence'' that can be established for them ([RRR] and Exp. IX); for example, we prove (Exp. IX) canonical isomorphisms such as
\[
K_\bullet(X[t])\simeq K_\bullet(X),\qquad \Gr_\bullet(X[t])\simeq \Gr_\bullet(X),
\]
which fail in general for $K^\bullet$ when $X$ is a nonregular scheme. On the other hand, $K^\bullet(X)$ and $\Gr^\bullet(X)$ lend themselves well to calculation because of their ring structures.

For schemes $X$ proper over a field, the projection
\[
f^X:X\longrightarrow(\operatorname{point})
\]
defines an ``\uline{Euler--Poincaré characteristic}'' homomorphism
\sgaStableTarget{sga6.e00.eq.6-4}\begin{equation}\tag{6.4}
\lambda_X=f^X_*:K_\bullet(X)\longrightarrow K_\bullet(\operatorname{point})=\mathbb Z,
\end{equation}
inducing on the subgroup $\Gr_0(X)=\operatorname{Filt}_0K_\bullet(X)$ the ``degree of $0$-cycles'' homomorphism
\sgaStableTarget{sga6.e00.eq.6-5}\begin{equation}\tag{6.5}
\deg_X:\Gr_0(X)\longrightarrow \mathbb Z.
\end{equation}
The multiplication homomorphism
\[
\Gr^i(X)\times\Gr_i(X)\longrightarrow \Gr_0(X),
\]
composed with the degree homomorphism (6.5), then gives a pairing
\sgaStableTarget{sga6.e00.eq.6-6}\begin{equation}\tag{6.6}
\Gr^i(X)\times \Gr_i(X)\longrightarrow \mathbb Z,
\end{equation}
which explains why, in certain respects, $\Gr^\bullet(X)$ and $\Gr_\bullet(X)$ play dual roles,

%%%%%%%%%%  scan idx 30  |  folio 17  %%%%%%%%%%
\SGAPrintPage{17}
just as cohomology and homology do. Thus, for a morphism $f:X\to Y$ of proper algebraic schemes, the two homomorphisms
\[
f^*:\Gr^\bullet(Y)\longrightarrow \Gr^\bullet(X)\qquad\text{and}\qquad f_*:\Gr_\bullet(X)\longrightarrow \Gr_\bullet(Y)
\]
are formally transposes of one another, with respect to the pairings (6.6).

The homomorphisms (6.4) and (6.5), and the pairing (6.6) obtained from them, can be regarded as the source of \uline{numerical relations} in intersection theory on proper algebraic schemes over a given field $k$. This point of view will be developed in some detail in Exp. X, which is largely independent of Exp. VI, VII, VIII (centred on the proof of the Riemann--Roch theorem) and of Exp. IX. Some notions concerning the \uline{behaviour of} $K^\bullet(X),\Gr^\bullet(X)$ \uline{under specialization} will also be given there.

\sgaStableTarget{sga6.e00.sec.7}\subsection*{7.}
As in every theory of intersections, it is important to specify the place of the \uline{Picard group} in the theory (or the group of \uline{divisor classes} in the classical terminology). Subject to a slight restriction, automatically satisfied for example when $X$ is quasi-compact, Exp. X establishes a canonical isomorphism
\sgaStableTarget{sga6.e00.eq.7-1}\begin{equation}\tag{7.1}
c_1:\Pic(X)\xrightarrow{\sim}\Gr^1(X),
\end{equation}
which shows the exact role of the Picard group in the theory we propose. This role justifies a more detailed study of this Picard group in Exposés XII, XIII, which are logically independent of the text of the Seminar. A study of \uline{numerical equivalence} and certain \uline{finiteness theorems for the Picard group} (announced without proof in [5]) is made by Kleiman in Exp. XIII, using techniques of a purely numerical nature (which do not use the theory of the Picard functor and its representability), due to himself and to Mumford, markedly simpler than Grothendieck's initial proofs. The only result refractory to Kleiman's methods (and which he is obliged to use) is [5, C-07 (i)], whose proof and certain refinements (including certain representability theorems for the Picard functor) occupy Exposé XII. The latter moreover has a spirit and a technique entirely foreign to the rest of the Seminar, being connected with [5]; this ex%
%%%%%%%%%%  scan idx 31  |  folio 18  %%%%%%%%%%
\SGAPrintPage{18}%
posé is logically independent of all the others.

Finally, in Exp. XI, of a distinctly more general nature than XII and XIII and closer to the general spirit of the Seminar, we study, on an arbitrary locally ringed topos, a remarkable functor called the ``\uline{determinant functor}'':
\sgaStableTarget{sga6.e00.eq.7-2}\begin{equation}\tag{7.2}
\det^*:\underline{\operatorname{Parfis}}(X)\longrightarrow \underline{\Inv}(X)
\end{equation}
from the category of perfect complexes on $X$, whose morphisms consist only of the isomorphisms in the sense of the derived category $D(X)$, to the category of invertible Modules on $X$. This functor is described almost completely by the requirement that it be of a ``local nature'' and that, for a bounded complex $L_\bullet$ with locally free components, it be given by the formula
\sgaStableTarget{sga6.e00.eq.7-3}\begin{equation}\tag{7.3}
\det^*(L_\bullet)\simeq\bigotimes_i \det(L_i)^{(-1)^i},
\end{equation}
where $\det(L_i)$ denotes the exterior power of maximal degree of the locally free Module $L_i$. Although this exposé too is logically independent of the rest of the Seminar except for Exposé I, it nevertheless forms an important part of the general ``yoga'' proposed in the Seminar and developed mainly in Exposés I to XI.

%%%%%%%%%%  scan idx 32  |  folio 19  %%%%%%%%%%
\SGAPrintPage{19}
\sgaStableTarget{sga6.e00.bibliography}\section*{\underline{BIBLIOGRAPHY}}
\begin{enumerate}[label={[\arabic*]},leftmargin=2.5em]
\item M. Atiyah and F. Hirzebruch, Vector bundles and homogeneous spaces, Symposia in Pure Mathematics, vol. 3, Differential Geometry p. 7--38, 1961.
\item A. Borel and J.P. Serre, Le théorème de Riemann--Roch, Bull. Soc. math. France, t. 86, p. 97--136 (1958).
\item H. Grauert, Ein Theorem der analytischen Garbentheorie, Pub. Math. no.~5, p. 5--64 (1960).
\item A. Grothendieck, Classes de faisceaux et théorème de Riemann--Roch, mimeographed notes, Princeton 1957, (cited as [RRR] and reproduced as an Appendix to Exp. 0 of the Seminar).
\item A. Grothendieck, Technique de descente et théorèmes d'existence en Géométrie Algébrique, in Fondements de la Géométrie Algébrique, Extracts from the Bourbaki Seminar 1957--1962, Secrétariat Mathématique, Paris.
\item W. Shih, Séminaire de Topologie à l'IHES, 1966/67.
\item M. Atiyah and F. Hirzebruch, The Riemann--Roch theorem for analytic embeddings, Topology, vol. 1, p. 151--166 (1962).
\end{enumerate}

\clearpage

